\title{Direct Preference Optimization: Your Language Model is Secretly a Reward Model}

\begin{abstract}
Large-scale unsupervised language models (LMs) acquire extensive world knowledge and exhibit some capacity for reasoning, yet exerting fine-grained control over their outputs remains challenging due to their wholly unsupervised training process. Conventional approaches to address this challenge involve gathering human annotations indicating relative preferences among model outputs, and subsequently fine-tuning the initial LM in line with these judgments—most commonly via reinforcement learning from human feedback (RLHF). Nevertheless, RLHF entails a multifaceted and frequently unstable pipeline: it first requires training a reward model that encapsulates human preferences, followed by using reinforcement learning to optimize the original LM for this estimated reward, all while preventing significant divergence from its pretrained state. This paper introduces an alternative reward model parameterization for RLHF, which facilitates the derivation of the optimal policy in closed form. This reparameterization reduces the RLHF framework to a straightforward classification objective. The proposed method, termed \textit{Direct Preference Optimization} (DPO), offers stability, strong empirical performance, and efficient computational requirements by removing the necessity for sampling from the LM during fine-tuning and minimizing reliance on extensive hyperparameter tuning. Empirical evaluations demonstrate that DPO is capable of fine-tuning LMs for human preference alignment on par with or superior to prevailing approaches. Remarkably, DPO outperforms PPO-based RLHF in managing the sentiment of generated text, and achieves comparable or better results in summarization and single-turn dialogue, while remaining notably simpler to implement and train.
\end{abstract}

\section{Introduction}
Recent advances in large unsupervised language models (LMs), trained on massive corpora, have demonstrated unexpected emergent abilities~\citep{chowdhery2022palm, brown2020language, touvron2023llama,bubeck2023sparks}. Nevertheless, the data underlying these models is sourced from humans who possess diverse intentions, skill levels, and objectives. As a result, the learned behaviors may reflect undesirable attributes; for instance, while it is beneficial for an AI coding tool to \textit{recognize} common programming errors so as to identify and correct them, we generally wish to steer code generation towards the subset of high-caliber, albeit infrequent, coding examples present in its training set. Similarly, even if it is important for a language model to \textit{know about} prevalent misconceptions held by, say, 50\% of people, we do not want it to propagate these inaccuracies in half the instances it encounters such queries. Consequently, the process of extracting the model’s \emph{preferred outputs and conduct} from its expansive reservoir of \textit{knowledge and competences} is a central challenge for ensuring AI systems that are secure, effective, and manageable \citep{ouyang2022training}. Common strategies rely on reinforcement learning (RL) to align LM behaviors with human values, but we demonstrate that the objectives used in current RL-based approaches can be exactly attained using a straightforward binary cross-entropy loss, thus streamlining the process of preference learning.

Broadly, standard approaches inject the targeted behaviors into a language model through carefully selected human preference datasets, emphasizing those qualities deemed safe and beneficial. This learning from preferences follows the initial, extensive phase of unsupervised pre-training on large-scale text collections. The simplest means of preference learning—supervised fine-tuning based on human-curated demonstrations of exemplary responses—has been surpassed in efficacy by reinforcement learning from human (or AI) feedback (RLHF/RLAIF; \citep{christiano2017deep,bai2022constitutional}). In RLHF frameworks, a reward model is fit to human preference judgments, after which RL is employed to adjust a policy so that it generates responses rated highly by the reward model, while remaining reasonably faithful to the original model distribution. Although RLHF yields models with notable strengths in conversation and code generation, its training regimen is substantially more intricate than that of standard supervised methods, requiring the training of multiple language models and integrating policy sampling during RL, thereby resulting in a high computational burden.

This work demonstrates a method for directly training language models to align with human preferences, bypassing the need for explicit reward modeling or reinforcement learning techniques. We introduce \textit{{Direct Preference Optimization} (DPO)}, a framework that, while implicitly maximizing the same objective as conventional RLHF approaches (reward maximization subject to a KL-divergence penalty), remains both straightforward to implement and efficient to train. The DPO update rule increases the log probability ratio between favored and unfavored outputs, but, crucially, incorporates a per-example, adaptive importance weight to mitigate the model collapse sometimes observed with simpler probability ratio methods. As with other approaches, DPO is grounded in a formal preference model (for instance, the Bradley-Terry model; \cite{bradley1952rankanalysis}), which provides a mechanism for quantifying the degree of agreement between a reward function and collected preference annotations. Nevertheless, where conventional methods leverage the preference model to guide reward model training and subsequently train policies to optimize the constructed reward, DPO exploits a change of variables to frame the preference loss directly in terms of the policy parameters. With access to human preference annotations over model outputs, DPO enables policy optimization via a binary cross entropy loss, thereby recovering an optimal policy relative to an implicit reward shaped by the preference data.

The primary contribution of this paper is the introduction of Direct Preference Optimization (DPO), an accessible, RL-free paradigm for preference-driven language model training. Empirical results indicate that DPO performs on par with, or surpasses, leading methods such as PPO-based RLHF, for tasks including sentiment control, summarization, and conversational modeling, across language models scaling up to 6B parameters.

\section{Related Work}

Self-supervised language models with ever-greater parameter counts have demonstrated the ability to perform certain tasks in a zero-shot setting \citep{radford2019language} or using few-shot prompts \citep{gpt3,megatron,chowdhery2022palm}. Nevertheless, enhancing their effectiveness on specific downstream applications and improving alignment with user goals can be achieved through further fine-tuning on curated datasets comprising instructions and human-generated completions \citep{mishra-etal-2022-cross,sanh2022multitask,chung2022scaling,thoppilan2022lamda}. This process, known as `instruction-tuning', equips LLMs to generalize beyond the distribution of instructions seen during fine-tuning, thus boosting their practical utility \citep{chung2022scaling}. Despite its successes, collecting \textit{relative} human evaluations of response quality typically proves more scalable than assembling datasets of expert-authored demonstrations. Consequently, later research has pursued fine-tuning LLMs with preference datasets to enhance performance in areas such as translation \citep{kreutzer-etal-2018-reliability}, summarization \citep{stiennon2022learning,ziegler2020finetuning}, narrative generation \citep{ziegler2020finetuning}, and following instructions \citep{ouyang2022training,ramamurthy2023is}. These approaches generally involve first training a neural reward model aligned with preference data—using a framework such as the Bradley-Terry model \citep{bradley1952rankanalysis}—and subsequently updating the language model parameters via reinforcement learning, with algorithms like REINFORCE \citep{williams1992reinforce}, proximal policy optimization (PPO; \cite{schulman2017proximal}), or their variants \citep{ramamurthy2023is}. Closely related strategies utilize LLMs—already fine-tuned to follow instructions and shaped by human feedback—to produce synthetic preference labels for properties such as safety or non-toxicity \citep{bai2022constitutional}. This process leverages minimal direct supervision, with humans providing only textual rubrics to guide the LLM's annotation process. Collectively, these techniques embody a synthesis of research traditions: one focused on reinforcement learning for optimizing language models toward diverse criteria~\citep{Ranzato2015SequenceLT,paulus2018a,wu2018learning}, and another on principled methodologies for learning from human preference data \citep{christiano2017deep,kupcsik2018learning}. Although leveraging relative human preferences offers notable advantages, applying reinforcement learning for fine-tuning large language models introduces significant practical obstacles; the present work contributes a theoretically motivated framework for optimizing with relative preferences absent explicit RL.

Outside the domain of language, the study of learning policies from preference information has been explored within both bandit frameworks and reinforcement learning paradigms, giving rise to a variety of methods. In the case of contextual bandit learning, where preferences or rankings of actions (as opposed to explicit reward signals) are available, the problem is formalized as the contextual dueling bandit (CDB; \cite{yue2012karmed,dudik2015contextual}). Without access to absolute reward values, theoretical work on CDBs replaces the standard notion of optimality with that of a \textit{von Neumann winner}—a policy guaranteed to achieve at least a 50\% win rate in expectation when compared against any alternative policy \citep{dudik2015contextual}. Nevertheless, a key distinction is that, in the CDB framework, preferences are received in an online fashion, whereas human preference learning usually relies on a static, offline collection of preference-labeled action pairs \citep{yan2022human}. Along similar lines, \textit{preference-based RL} (PbRL) utilizes binary preferences inferred from an unspecified ‘scoring’ function, rather than relying on reward signals \citep{BusaFekete2014,ruiz2023dueling}. A range of PbRL algorithms have been developed, with many supporting the use of off-policy preference data; however, they generally involve two stages: first modeling the latent scoring or reward function and then performing optimization based on it \citep{jain2013learning,BusaFekete2014,christiano2017deep,sadigh2017active,kupcsik2018learning}. In contrast, we propose a unified, single-step procedure for policy optimization that aims to directly maximize alignment with observed preferences.

\section{Preliminaries}\label{section:prelims}

We summarize the RLHF process as described by \citeauthor{ziegler2020finetuning} and subsequently by others \citep{stiennon2022learning, bai2022training, ouyang2022training}. This process generally consists of three distinct stages: (1) supervised fine-tuning (SFT); (2) the collection of preference data along with reward model training; and (3) reinforcement learning-based policy optimization.

\textbf{SFT}: The RLHF workflow typically commences with adapting a pre-trained language model to a particular downstream application (such as conversation or summarization) via supervised learning, producing a model denoted $\pi^\text{SFT}$.

\textbf{Reward Modelling Phase}: During the second stage, the SFT-tuned model receives input prompts $x$ and generates candidate response pairs $(y_1, y_2)\sim \pi^\text{SFT}(y \mid x)$. These outputs are subsequently shown to human evaluators, who select their preferred answer (notated as $y_w\succ y_l \mid x$, with $y_w$ and $y_l$ denoting the chosen and non-chosen responses, respectively). It is assumed that these preference labels are outcomes of some underlying, unknown reward model $r^*(y, x)$. Various models for encoding preferences have been employed, with the Bradley-Terry (BT) model \cite{bradley1952rankanalysis} being widely adopted; more broadly, Plackett-Luce ranking models \citep{plackett1975analysis, luce2012individual} are also applicable should multiple candidate rankings be available. The BT framework specifies that the true human preference probabilities $p^*$ can be expressed as follows:

Given a fixed dataset of comparison pairs $\mathcal{D}=\bigl\{x^{(i)}, y_w^{(i)}, y_l^{(i)}\bigr\}_{i=1}^N$ sampled from $p^*$, a reward model $r_{\phi}(x, y)$ can be specified, with its parameters estimated through the maximum likelihood principle. When posed as a binary classification problem, the associated negative log-likelihood loss can be written as:

\begin{equation}\label{eq:reward_model}
    \mathcal{L}_R(r_{\phi}, \mathcal{D}) = -\mathbb{E}_{(x, y_w, y_l)\sim \mathcal{D}}\bigl[\log \sigma(r_{\phi}(x, y_w)- r_{\phi}(x, y_l))\bigr]
\end{equation}

where $\sigma$ is the logistic sigmoid function. Within large language models, the architecture for $r_{\phi}(x, y)$ is usually built upon the SFT model $\pi^\text{SFT}(y \mid x)$, to which a final linear projection is appended atop the output of the last transformer block, generating a scalar reward value~\cite{ziegler2020finetuning}. To reduce variance in the resulting reward scores, it is standard to normalize rewards so that $\mathbb{E}_{x,y\sim \mathcal{D}}\left[r_\phi(x, y)\right] = 0$ for each $x$.

\textbf{RL Fine-Tuning Phase}: In the reinforcement learning phase, the trained reward model guides updates to the language model’s parameters. Consistent with prior works~\citep{jaques2017sequence, jaques2020human}, the objective takes the following form:

\begin{equation}\label{eq:RL}
\max_{\pi_{\theta}}  \mathbb{E}_{x\sim \mathcal{D}, y\sim \pi_{\theta}(y \mid x)}\bigl[r_{\phi}(x, y)\bigr] - \beta\mathbb{D}_{\textrm{KL}}\bigl[\pi_{\theta}(y\mid x)\mid \mid \pi_\text{ref}(y\mid x)\bigr],
\end{equation}

where the parameter $\beta$ mediates the trade-off with respect to the KL-divergence from a reference policy $\pi_\text{ref}$, generally chosen as the original SFT model $\pi^\text{SFT}$. Practically, $\pi_\theta$ is also initialized from $\pi^\text{SFT}$. The KL regularization is essential, discouraging the policy from straying too far from regions where the reward model remains reliable, and helping to retain output diversity and to avoid collapse to a single highly-rewarded response. Owing to the discrete sampling inherent in language generation, this objective function is not amenable to gradient-based optimization directly, and reinforcement learning approaches are employed. The conventional method~\citep{ziegler2020finetuning, stiennon2022learning, bai2022training, ouyang2022training} formulates the reward as ${r(x, y) = r_{\phi}(x, y) -\beta (\log \pi_{\theta}(y\mid x) - \log \pi_\text{ref}(y\mid x))}$, which is optimized via PPO~\cite{schulman2017proximal}. 

\section{Direct Preference Optimization}\label{sec:DPO}

Recognizing the practical barriers that arise in scaling RL-based methods for language model fine-tuning, we set out to formulate a straightforward method for preference-based policy optimization. Rather than following conventional RLHF paradigms—where a reward model is learned and subsequently optimized via RL—we instead propose a specific parameterization of the reward that allows us to directly characterize the optimal policy in closed-form, eliminating the need for a full RL training loop. Our central insight is to use an analytic connection between the reward function and its corresponding optimal policy; this allows us to translate a loss over the reward model into a loss over the policy space. By adopting this change of variables, we bypass the need for an explicit reward model, yet our approach still conforms to standard preference models such as Bradley-Terry. As a result, our policy network effectively encompasses both the generative language model and its underlying (implicit) reward structure.

\textbf{Derivation of the DPO objective.} As an initial step, we consider the RL objective presented in previous studies, specifically Eq.~\ref{eq:RL}, while adopting a general reward function $r$. Consistent with the methodologies of earlier works~\citep{peters2007reinforcement, peng2019advantage, korbak2022reinforcement, go2023aligning}, it is a direct result that the KL-regularized reward maximization problem defined in Eq.~\ref{eq:RL} admits an optimal policy expressed as:

\begin{equation}\label{eq:op_policy}
    \pi_r(y\mid x) = \frac{1}{Z(x)}\pi_\text{ref}(y\mid x)\exp\left(\frac{1}{\beta}r(x, y)\right),
\end{equation}

where the normalizing constant $Z(x) = \sum_{y}\pi_\text{ref}(y\mid x)\exp\left(\frac{1}{\beta}r(x, y)\right)$ is commonly referred to as the partition function. For a comprehensive step-by-step derivation, consult Appendix \ref{app:derivation1}. Even assuming availability of an MLE-based estimator $r_{\phi}$ for the true reward function $r^*$, direct computation of $Z(x)$ remains computationally prohibitive \citep{korbak2022reinforcement, go2023aligning}, thereby limiting the practicality of this policy formulation. Nonetheless, we can algebraically manipulate Eq.~\ref{eq:op_policy} to express the reward function as a function of its optimal policy $\pi_r$, the baseline $\pi_\text{ref}$, and $Z(\cdot)$. In particular, by applying a logarithmic transformation and rearranging, we arrive at:

\begin{equation}\label{eq:main_eq}
    r(x,y) =\beta \log \frac{\pi_r(y\mid x)}{\pi_\text{ref}(y\mid x)} + \beta \log Z(x).
\end{equation}

This representation can be instantiated for the ground-truth reward $r^*$ and its associated optimal policy $\pi^*$. Importantly, under the Bradley-Terry model, only differences between reward values for two candidate completions matter, that is, ${p^*(y_1 \succ y_2 \mid x) = \sigma(r^*(x, y_1) - r^*(x, y_2))}$. Plugging in the parametrization from Eq.~\ref{eq:main_eq} for $r^*(x,y)$ in the preference probability given by Eq.~\ref{eq:bradley-terry}, the term involving the partition function $Z(x)$ drops out, yielding a preference model defined purely in terms of $\pi^*$ and the reference policy $\pi_\text{ref}$. Thus, under the Bradley-Terry construction, the optimal policy $\pi^*$ for RLHF adheres to:

\begin{equation}\label{eq:objective}
    p^*(y_1\succ y_2 \mid x)=\frac{1}{1 + \exp\left(\beta \log \frac{\pi^*(y_2\mid x)}{\pi_\text{ref}(y_2\mid x)} - \beta \log \frac{\pi^*(y_1\mid x)}{\pi_\text{ref}(y_1\mid x)}\right)}
\end{equation}

The corresponding mathematical development is detailed in Appendix~\ref{app:derivation2}. Although Eq.~\ref{eq:objective} specifically relies on the Bradley-Terry model, an analogous treatment applies for more general preference frameworks such as the Plackett-Luce model~\citep{plackett1975analysis, luce2012individual}; see Appendix~\ref{app:plackett_luce_models} for further elaboration.

Having cast the preference likelihood in terms of the policy instead of the reward function, we can define a maximum likelihood estimator for any parameterized policy $\pi_\theta$. Mirroring the maximum likelihood formulation for reward models (as in Eq.~\ref{eq:reward_model}), this leads to the following policy objective:

\begin{equation}\label{eq:optimum_model}
    \mathcal{L}_\text{DPO}(\pi_{\theta}; \pi_\text{ref}) = -\mathbb{E}_{(x, y_w, y_l)\sim \mathcal{D}}\left[\log \sigma \left(\beta \log \frac{\pi_{\theta}(y_w\mid x)}{\pi_\text

Through this approach, we learn an implicit reward via a different parameterization, where the resulting optimal policy corresponds directly to $\pi_\theta$. Notably, the process aligns with learning a reparameterized Bradley-Terry model, thus inheriting well-established theoretical guarantees, such as consistency under appropriate preference data assumptions \cite{bong2022generalized}. Further discussion of DPO's theoretical aspects and its connections to related literature can be found in Section~\ref{sec:theory}.

\textbf{Mechanistic view of the DPO update.} To better grasp how DPO functions, we examine the gradient of the loss function $\mathcal{L}_\text{DPO}$. The gradient with respect to the model parameters $\theta$ is given by:

\begin{multline*}\label{eq:gradient}
    \nabla_\theta \mathcal{L}_\text{DPO}(\pi_\theta;\pi_\text{ref}) = \\ -\beta\mathbb{E}_{(x, y_w, y_l) \sim \mathcal{D}} \bigg[\underbrace{\sigma(\hat{r}_\theta(x, y_l) - \hat{r}_\theta (x, y_w))}_\text{emphasis increases if the reward model's ordering is incorrect}\bigg[\underbrace{\nabla_\theta\log \pi(y_w \mid x)}_\text{encourage preferred completion $y_w$} - \underbrace{\nabla_\theta\log\pi(y_l \mid x)}_\text{discourage dispreferred completion $y_l$}\bigg]\bigg],
\end{multline*}

where the implicit reward, $\hat{r}_\theta(x, y) = \beta \log \frac{\pi_\theta(y \mid x)}{\pi_\text{ref}(y \mid x)}$, is determined by the language model $\pi_\theta$ relative to the reference policy $\pi_\text{ref}$ (details provided in Section~\ref{sec:theory}). In essence, this gradient update pushes up the probability of preferred continuations $y_w$ and reduces it for less favored ones $y_l$. The update is modulated by how much more highly the implicit reward function ranks the non-preferred responses, and this sensitivity is amplified by $\beta$, reflecting both model misorderings and the KL constraint's influence. Empirically, this adaptive weighting proves crucial: omitting it—i.e., removing the coefficient—often leads to problematic model behaviors (see Appendix Table~\ref{tab:unlikelihood_generations}).

\textbf{DPO overview.} The standard DPO workflow comprises the following steps: (1) For each prompt $x$, generate candidate completions $y_1, y_2 \sim \pi_\text{ref}(\cdot \mid x)$ and annotate them according to human preferences to create an offline preference dataset $\mathcal{D} = \{x^{(i)}, y_w^{(i)}, y_l^{(i)}\}_{i=1}^N$. (2) Adjust the language model $\pi_\theta$ by minimizing the DPO loss $\mathcal{L}_\text{DPO}$, conditioned on $\pi_\text{ref}$, the assembled dataset $\mathcal{D}$, and a specified $\beta$ value. In practical settings, there is a preference for utilizing publicly available preference datasets rather than collecting new human feedback and completions. Because such datasets are typically obtained using $\pi^\text{SFT}$, we set $\pi_\text{ref} = \pi^\text{SFT}$ whenever feasible. If $\pi^\text{SFT}$ is inaccessible, we initialize $\pi_\text{ref}$ by maximizing the likelihood of preferred completions ${(x, y_w)}$, that is, by solving ${\pi_\text{ref} = \argmax_{\pi}\mathbb{E}_{x, y_w \sim \mathcal{D}}\left[\log \pi(y_w \mid x)\right]}$. This approach reduces mismatch between the inaccessible true reference distribution and the surrogate $\pi_\text{ref}$ employed in DPO. Supplementary information regarding implementation specifics and hyperparameter choices can be found in Appendix~\ref{app:implementation}.

\section{Theoretical Analysis of DPO}
In this section, we elaborate on the interpretation of DPO, present theoretical justifications, and connect the benefits of DPO to shortcomings in actor-critic methods employed for RLHF (including PPO~\cite{schulman2017proximal}).

\label{sec:theory}

\subsection{Your Language Model Is Secretly a Reward Model} DPO replaces the explicit reward modeling and reinforcement learning phases typically required for policy training with a unified maximum likelihood objective. It is important to observe that the loss function defined in Eq. \ref{eq:main_eq} matches the structure of a Bradley-Terry model under the reward transformation $r^*(x, y) = \beta \log\frac{\pi^*_\theta(y \mid x)}{\pi_\text{ref}(y \mid x)}$. Here, the parametric model $\pi_{\theta}$ is optimized analogously to the reward model optimization outlined in Eq. \ref{eq:reward_model} via a change of variables. In what follows, we develop the theoretical framework underlying this reparameterization, demonstrate that it imposes no additional constraints on the family of reward models that can be learned, and prove that this process enables exact retrieval of the optimal policy. We initiate the analysis by formalizing an equivalence relation over reward functions.

\begin{definition}
Two reward functions $r(x, y)$ and $r'(x, y)$ are considered equivalent if and only if there exists a function $f$ such that $r(x, y)-r'(x, y) = f(x)$.
\end{definition}

This equivalence defines a partitioning of the space of reward functions, dividing them into equivalence classes. We present the following two lemmas:

\begin{lemma}\label{lemma:same_prefrence}
Within the Plackett-Luce, including the Bradley-Terry, preference model, any two reward functions from the same equivalence class will yield identical preference distributions.
\end{lemma}

\begin{lemma}\label{lemma:same_policy}
    Any two reward functions belonging to the same equivalence class will result in the same optimal policy for the corresponding constrained RL problem.
\end{lemma}

The proofs, which follow directly from the definitions, are presented in Appendix \ref{app:lemma1}. The first lemma corresponds to a classical under-identification issue observed with the Plackett-Luce class of models \cite{plackett1975analysis}. Owing to this ambiguity, additional identifiability conditions are commonly imposed in order to ensure any guarantees concerning the MLE solutions derived from Eq. \ref{eq:reward_model} \cite{bong2022generalized}. The second lemma implies that every reward function within the same equivalence class induces the same optimal policy, so when targeting policy recovery, it suffices to recover any representative reward function from the optimal equivalence class. The following theorem is established in Appendix~\ref{app:thm1}:

\begin{theorem}\label{thm:main}
    Assuming mild conditions, all reward equivalence classes compatible with the Plackett-Luce (and in particular, the Bradley-Terry) frameworks can be represented through the parameterization ${r(x, y) = \beta \log \frac{\pi(y\mid x)}{\pi_\text{ref}(y\mid x)}}$ for some choice of model $\pi(y\mid x)$ and reference model $\pi_\text{ref}(y \mid x)$.
\end{theorem}

\begin{sproof}
    Let $r(x, y)$ be an arbitrary reward function inducing an optimal policy $\pi_r(y \mid x)$, as defined by Eq. \ref{eq:op_policy}. We demonstrate that a representative reward function from its equivalence class admits the form specified above. Define the mapping $f$ as follows:
\begin{equation}
    f(r; \pi_\text{ref}, \beta)(x, y) = r(x, y) - \beta\log\sum_{y}\pi_\text{ref}(y\mid x)\exp\left(\frac{1}{\beta}r(x, y)\right)
\end{equation}
Here, $f$ adjusts the original reward by subtracting the logarithm of the partition function (dependent only on $x$), normalizing $r$. Because this normalization is solely a function of $x$, $f(r; \pi_\text{ref}, \beta)(x, y)$ remains in the same equivalence class as $r(x, y)$. If we now set $r$ according to the right-hand side of Eq.~\ref{eq:main_eq}, which is valid for any reward function, we obtain $f(r; \pi_\text{ref}, \beta)(x, y) = \beta \log \frac{\pi_r(y\mid x)}{\pi_\text{ref}(y\mid x)}$. Thus, the mapping $f$ yields a representative in the equivalence class of $r$ exhibiting the proposed parameterization, ensuring that this reparameterization does not restrict the generality of the reward model.
\end{sproof}

An alternative interpretation of Theorem~\ref{thm:main} is that it delineates precisely which representative of a given equivalence class of reward functions is chosen via the DPO reparameterization; specifically, it is the reward function satisfying

\begin{equation}\label{eq:lag_p}
     \sum_{y}\underbrace{\pi_\text{ref}(y\mid x)\exp\left(\frac{1}{\beta}r(x, y)\right)}_{=\pi(y\mid x)\text{, using Thm.~\ref{thm:main} reparam.}} = 1,
\end{equation}

meaning that $\pi(y\mid x)$ constitutes a well-defined probability distribution (with non-negative probabilities summing to unity). Furthermore, referring to Eq.~\ref{eq:op_policy}, it is evident that Eq.~\ref{eq:lag_p} defines the partition function associated with the optimal policy corresponding to the reward function $r(x, y)$. The central observation underlying the DPO algorithm is that, by imposing suitable constraints on the otherwise under-specified Plackett-Luce (notably, the Bradley-Terry) class of preference models, one can retain the expressive power regarding possible reward models, while concurrently ensuring that the optimal policy described in Eq.~\ref{eq:op_policy} is computationally tractable for any prompt $x$.

\subsection{Instability of Actor-Critic Algorithms}
Our theoretical framework also facilitates an analysis of the instabilities exhibited by conventional actor-critic approaches within RLHF, including PPO. We focus on the RL fine-tuning phase as detailed in Section~\ref{section:prelims}, and relate this to the control as inference paradigm \cite{levine2018reinforcement} for the constrained RL problem presented in \ref{eq:RL}. Considering a parameterized policy $\pi_{\theta}(y\mid x)$, we seek to minimize $\mathbb{D}_{\text{KL}}[\pi_{\theta}(y|x) \mid \mid \pi^*(y\mid x)]$, where $\pi^*$ represents the optimal policy arising from Eq.~\ref{eq:optimum_model} with reward function $r_{\phi}(y, x)$. Through algebraic manipulation, this yields the following objective for optimization:

\begin{equation}\label{eq:AC}
    \max_{\pi_{\theta}}\mathbb{E}_{\pi_{\theta}(y\mid x)}\bigg[\underbrace{r_{\phi}(x, y) -\beta\log\sum_{y}\pi_\text{ref}(y\mid x)\exp\left(\frac{1}{\beta}r_{\phi}(x, y)\right)}_{f(r_{\phi}, \pi_\text{ref}, \beta)} - \underbrace{\beta\log\frac{\pi_{\theta}(y\mid x)}{\pi_\text{ref}(y\mid x)}}_{\text{KL}}\bigg]
\end{equation}

This objective matches that of previous studies 
\citep{ziegler2020finetuning, stiennon2022learning, bai2022training, ouyang2022training}, which optimize with the DPO-equivalent reward formulated for the reward class $r_{\phi}$. In this context, the normalization factor within $f(r_{\phi}, \pi_\text{ref}, \beta)$ can be viewed as representing the soft value function associated with the reference policy $\pi_\text{ref}$. Although the normalization does not influence the final optimum, omitting it can introduce significant variance into the policy gradient estimates, thereby destabilizing the learning process. One possible remedy is to use an estimated value function for normalization, though this approach presents its own optimization challenges. Previous research has instead normalized rewards by employing a human completion as a baseline—effectively using a Monte Carlo estimate with a single sample for the normalizing constant. Distinctly, the DPO reparameterization produces a reward function inherently free from reliance on any such baseline.

\section{Experiments}
Here, we conduct empirical studies to assess the effectiveness of DPO in learning policies directly from preference data. To begin, in a controlled text generation environment, we investigate how DPO balances reward maximization against maintaining low KL-divergence with the reference policy, comparing its efficiency with established preference learning methods like PPO. We then test DPO on larger models and more challenging RLHF benchmarks, including tasks such as summarization and dialogue. Remarkably, we observe that with minimal hyperparameter adjustment, DPO consistently matches or surpasses the performance of robust baselines, such as RLHF using PPO, and outperforms methods that select the best trajectory from $N$ samples based on a learned reward model. Prior to presenting our main results, we outline the experimental framework; additional specifics are included in Appendix~\ref{app:exp_details}.

\textbf{Tasks.} We evaluate our methods on three distinct open-ended text generation tasks. Across all tasks, the learning algorithms derive a policy based on a preference dataset, denoted by $\mathcal{D}=\bigl\{x^{(i)}, y_w^{(i)}, y_l^{(i)}\bigr\}_{i=1}^N$. In the task of \textbf{controlled sentiment generation}, the input $x$ serves as a prompt extracted from an IMDb movie review \cite{maas-EtAl:2011:ACL-HLT2011}, with the objective for the policy to produce a continuation $y$ exhibiting positive sentiment. To enable a controlled assessment, we synthesize preference pairs by generating outputs and comparing them using a pre-trained sentiment classifier such that $p(\text{positive}\mid x,y_w)>p(\text{positive}\mid x,y_l)$. For supervised fine-tuning (SFT), we adapt GPT-2-large on the training subset of the IMDb corpus until performance plateaus (implementation specifics appear in App~\ref{app:sentiment_details}). In the \textbf{summarization} task, $x$ corresponds to a Reddit forum post, and the policy is tasked with generating a summary $y$ capturing the primary arguments. Consistent with earlier research, we employ the Reddit TL;DR summarization corpus \citep{volske-etal-2017-tl} augmented with preference judgments collected by \citeauthor{stiennon2022learning}. Our SFT baseline leverages a model trained on human-authored Reddit post summaries\footnote{\url{https://huggingface.co/CarperAI/openai_summarize_tldr_sft}} and is optimized using the TRLX \citep{leandro_von_werra_2023_7790115} library for RLHF. Notably, the preference labels used stem from a dataset assembled by \citeauthor{stiennon2022learning}, involving samples produced by a closely-related, albeit distinct, SFT model. In the case of \textbf{single-turn dialogue}, each $x$ represents a user prompt—ranging from scientific queries to advice requests—and the goal for the policy is to generate a reply $y$ that is both engaging and supportive. We utilize the Anthropic Helpful and Harmless dialogue dataset \citep{bai2022training}, which features 170k interactions between humans and an AI assistant; every conversation terminates with two alternative assistant replies, along with an annotation indicating human preference. Since no pre-existing SFT model is available in this domain, we train one by fine-tuning an existing language model exclusively on the human-preferred responses.

\textbf{Evaluation.} We employ two distinct strategies to assess performance in our experiments. To gauge how effectively each algorithm meets the goal of constrained reward maximization, we use a controlled sentiment generation framework where we compute the Pareto frontier of reward obtained versus KL-divergence from the reference policy; this calculation is enabled by access to the ground-truth reward function, implemented as a sentiment classifier. In contrast, since the actual reward function is typically unavailable in practical applications, we assess the algorithms by measuring their \textit{win rate} when compared to a baseline policy, utilizing GPT-4 as a stand-in for human evaluators in the contexts of summary quality (summarization) and helpfulness (single-turn dialogue). The baseline for summarization is provided by the test set’s reference summaries, and for dialogue by the response preferred in the test data. Although previous research has indicated that LMs may serve as more reliable automatic evaluators than existing automated metrics \citep{Chen2023ExploringTU}, we further support our adoption of GPT-4 for evaluation through a human study, described in Sec.~\ref{sec:human-judgments}. Our findings demonstrate a high degree of correlation between GPT-4 and human judgments, with agreement levels between humans and GPT-4 equaling or exceeding those observed among human annotators.

\textbf{Methods.} Beyond DPO, we compare a range of established methods for tuning language models in alignment with human preferences. As a simple baseline, we utilize zero-shot prompting with \textbf{GPT-J} \citep{gpt-j} for summarization tasks, and 2-shot prompting with \textbf{Pythia-2.8B} \citep{biderman2023pythia} for dialogue generation. We also assess the \textbf{SFT} approach and introduce \textbf{Preferred-FT}, which involves supervised fine-tuning on the selected completion $y_w$, sourced either from SFT (in sentiment and summarization scenarios) or from a generic language model (in single-turn dialogue). An additional pseudo-supervised approach is \textbf{Unlikelihood}~\citep{welleck2019neural}, which is trained to both increase the likelihood of $y_w$ and decrease that of $y_l$, with an optional coefficient $\alpha\in[0,1]$ applied to the “unlikelihood” component. We further incorporate \textbf{PPO} \citep{schulman2017proximal}, employing a reward model trained on preference data, as well as \textbf{PPO-GT}, which serves as an oracle utilizing access to the ground-truth reward function in the sentiment-controlled scenario. For the sentiment experiments, two variants of PPO-GT are applied: an off-the-shelf implementation \cite{leandro_von_werra_2023_7790115} and a modified version featuring normalized rewards and enhanced hyperparameter tuning, adjustments which are also adopted when running PPO with learned rewards. Lastly, we evaluate the \textbf{Best of $N$} baseline, where $N$ outputs are generated from either the SFT or Preferred-FT models (for dialogue), and the response with the highest score under a preference-trained reward model is selected. While this technique can deliver superior results, it is computationally intensive since it necessitates $N$ completions for every input at inference time, even with moderate values of $N$.

\subsection{How well can DPO optimize the RLHF objective?}

Typical RLHF algorithms employ a KL-constrained reward maximization framework, which encourages the policy to maximize reward while maintaining proximity to the reference policy. Thus, when evaluating and comparing these algorithms, both the obtained reward and the extent of KL divergence are important: achieving marginally better rewards at the cost of a substantially larger KL is often suboptimal. In Figure~\ref{fig:frontier-tldr-main}, the tradeoff between reward and KL—the reward-KL frontier—is presented for several algorithms within the sentiment classification context. For each method, we conduct several independent training sessions, varying the policy conservativeness parameter for each run (for PPO, target KL $\in\{3,6,9,12\}$; for DPO, $\beta \in \{0.05,0.1,1,5\}$; for unlikelihood, $\alpha\in\{0.05,0.1,0.5,1\}$; and random seeds for preferred-FT), totaling 22 experiments. Policies are periodically evaluated every 100 training iterations on a dedicated set of test prompts, measuring both the mean reward under the true reward function and the mean sequence-level KL\footnote{Calculated as the aggregate per-timestep KL-divergence across sequences.} with respect to the reference policy, $\text{KL}\left(\pi\mid \mid \pi_\text{ref}\right)$. Our analysis reveals that DPO achieves a reward-KL frontier that is markedly more efficient than other approaches, reaching superior rewards with comparatively low KL divergence. This observation is striking for several reasons. Notably, while DPO and PPO share an identical objective, DPO demonstrates far greater efficiency, strictly outperforming PPO in the reward versus KL tradeoff. Furthermore, DPO surpasses PPO even in settings where PPO is provided with access to ground truth rewards (PPO-GT).

\subsection{Can DPO scale to real preference datasets?}
\label{sec:dpo-real-datasets}
We proceed to assess how well DPO performs when fine-tuned on real-world preference datasets in tasks like summarization and single-turn dialogue. For summarization, traditional automatic metrics such as ROUGE are known to align weakly with human judgment~\citep{stiennon2022learning}. Earlier studies have demonstrated that fine-tuning language models with PPO using human feedback can yield more compelling summaries. In our evaluation, we generate outputs on the test partition of the TL;DR summarization dataset and calculate the average rate at which these outputs are favored over reference completions. For all approaches, completions are drawn at sampling temperatures between 0.0 and 1.0, with the resulting win rates depicted in Figure~\ref{fig:frontier-tldr-main} (right). DPO, PPO, and Preferred-FT all use the same GPT-J SFT model as their initialization point\footnote{\url{https://huggingface.co/CarperAI/openai_summarize_tldr_sft}}. Our results indicate that DPO achieves a win rate close to 61\% at a sampling temperature of 0.0, surpassing PPO's peak win rate of approximately 57\% at the same temperature. Additionally, DPO attains a greater maximum win rate compared to the best-of-$N$ reference baseline. It should be noted that we did not undertake extensive tuning of the $\beta$ hyperparameter for DPO, which suggests that its performance could potentially be further improved. Notably, DPO also displays greater resilience to variations in sampling temperature than PPO, whose effectiveness diminishes towards the baseline GPT-J performance at higher temperatures. In contrast, Preferred-FT offers minimal gains over the SFT baseline. Furthermore, as described in Section~\ref{sec:human-judgments}, direct human comparison reveals that DPO samples at temperature 0.25 are chosen 58\% of the time over those from PPO at the same temperature.

For the single-turn dialogue task, we assess the various approaches on the segment of the Anthropic HH dataset \citep{bai2022training} test split that consists of one exchange between human and assistant. For GPT-4 assessments, we determine win rates by comparing model outputs to the preferred completions in the test set, which serve as ground truth. Due to the lack of an established SFT model for this particular task, our procedure involves initializing from a pre-trained Pythia-2.8B, fine-tuning a reference model using Preferred-FT on the preferred responses to ensure sampled completions match the reference distribution, and subsequently applying DPO training. We further benchmark our results against the strongest out of 128 Preferred-FT completions (as performance for the Best of $N$ baseline was observed to plateau at $N=128$; consult Appendix Figure~\ref{fig:best-of-n}), as well as a 2-shot prompt setup with the Pythia-2.8B base model. DPO demonstrates performance equal to or superior to the other approaches at their respective optimal sampling temperatures. Additionally, we evaluate an RLHF model using PPO trained on the Anthropic HH dataset\footnote{\url{https://huggingface.co/reciprocate/ppo_hh_pythia-6B}} sourced from a widely recognized repository\footnote{\url{https://github.com/CarperAI/trlx/tree/main/examples/hh}}, but are unable to identify any prompt or sampling temperature that surpasses the performance of the base Pythia-2.8B. Given our findings with TL;DR and recognizing that both methods aim to maximize the same reward function, we treat the Best of 128 baseline as an approximate stand-in for PPO-level results. In summary, DPO stands out as the only computationally efficient technique to yield improvement over preferred completions from the Anthropic HH dataset, while matching or exceeding the performance of the computationally heavy Best of 128 strategy. Lastly, Figure~\ref{fig:dialogue-main} indicates that DPO rapidly reaches optimal performance.

\subsection{Generalization to a new input distribution}

To deepen our analysis of how PPO and DPO perform when faced with distributional changes, we assess the policies derived from our Reddit TL;DR summarization task on an alternate domain: news articles found in the test partition of the CNN/DailyMail dataset \citep{nallapati-etal-2016-abstractive}. For this evaluation, we employ the optimal sampling temperatures determined from the TL;DR experiment (0 and 0.25). The findings, displayed in Table~\ref{tab:ood}, indicate that DPO policies continue to significantly outperform PPO under these shifted conditions. To measure performance, we use the same GPT-4 (C) prompt as in the TL;DR setting, modifying “forum post” to “news article,” and compute the rate at which GPT-4 preferred model summaries over ground-truth references. These results offer preliminary support that DPO policies can generalize at least as effectively as PPO policies, despite DPO not relying on the supplementary unlabeled Reddit TL;DR prompts used by PPO.

\subsection{Validating GPT-4 judgments with human judgments}
\label{sec:human-judgments}
We run a human evaluation to test how reliably GPT-4’s preferences align with those of humans, drawing on outcomes from the TL;DR summarization study and leveraging two distinct prompts for GPT-4. The first, \textbf{GPT-4 (S)} (simple), asks which summary better covers the essential information of the post. The second, \textbf{GPT-4 (C)} (concise), extends this by asking for the more concise summary—a criterion included since we observed that GPT-4 tends to prefer lengthier, more redundant summaries than humans when using only the simple prompt. Full prompt texts are provided in Appendix~\ref{app:prompts}. We execute three comparative evaluations using: the top-performing method (DPO, temp. 0.25), the lowest (PPO, temp. 1.0), and a method with intermediate quality (SFT, temp. 0.25), each measured against PPO sampled greedily (its optimal temperature). Our findings indicate that, under both prompting strategies, GPT-4’s preferences are as consistent with human judgments as inter-human agreement, supporting GPT-4's suitability as a stand-in for human assessors. Given limited human annotator resources, multiple judgments were gathered for only the DPO and PPO-1 comparisons. Notably, the \textbf{GPT-4 (C)} prompt usually produces win rates that more closely reflect human evaluations, and thus we adopt it for reporting the principal results in Section~\ref{sec:dpo-real-datasets}. For expanded information regarding our human evaluation process, the rating interface, and the roster of annotators, see Appendix~\ref{app:human-study}.

\section{Discussion}
Preference-based learning presents an effective and scalable approach for developing robust and aligned language models. In this work, we introduced DPO, a straightforward methodology for leveraging human preferences in language model training without resorting to reinforcement learning techniques. Rather than adapting the preference learning task to fit the conventional RL paradigm to utilize existing RL frameworks, DPO proposes an explicit correspondence between policy distributions and reward functions. This enables direct optimization of human-aligned behaviors in language models through a simple cross-entropy loss, completely bypassing reinforcement learning while preserving the problem’s expressiveness. Remarkably, with minimal adjustment of hyperparameters, DPO matches or surpasses the performance of traditional RLHF algorithms, such as those employing PPO, effectively lowering the obstacles to preference-driven language model training.

\textbf{Limitations \& Future Work.} The findings in this paper open several avenues for further investigation. A critical question is how well DPO-trained policies extrapolate to unseen data relative to those trained via explicit reward modeling. Preliminary evidence points to generalization abilities on par with PPO-trained models, yet a thorough analysis remains necessary. For instance, it is worth investigating whether leveraging self-generated labels from the DPO policy can enhance the use of unlabeled prompts. Another question pertains to the dynamics of reward over-optimization in the context of direct preference optimization and whether the moderate performance dip observed in Figure~\ref{fig:dialogue-main}-right is indicative of such effects. Moreover, our experiments focus on models up to 6B parameters; expanding DPO to much larger, cutting-edge models constitutes a promising research trajectory. Regarding assessment methodologies, we noticed that GPT-4’s computed win rates are prompt-dependent; future research might address optimal strategies for extracting reliable judgments from automated evaluators. Ultimately, the scope of DPO may extend beyond human preference training for language models, offering potential benefits for generative modeling across various domains.